\documentclass[11pt,a4paper]{article}

% --- ENCODING & LANG ---
\usepackage[utf8]{inputenc}
\usepackage[T1]{fontenc}
\usepackage[ngerman,main=english]{babel}
\usepackage{lmodern}
\usepackage{microtype}

% --- MATH & GRAPHICS ---
\usepackage{amsmath,amssymb,amsfonts}
\usepackage{graphicx}
\usepackage{booktabs}
\usepackage{cite}
\usepackage{hyperref}
\numberwithin{table}{section}
\usepackage{float}

\title{\textbf{Nicht-Metaphysics: Paper I - The Apophatic Riemann Invariant}\\[0.3em]
\large Non-Riemannian Counterexamples, Tautological Friction Metrics,\\
and Phase-Coupled Spectral Rigidity}
\author{\textbf{xerx593} \quad \& \quad \textbf{Non-Human Interlocutors}}
\IfFileExists{release_dates.def}{% ======================================================================
% Central Production Release Date Registry for nicht-physics
% Maps root filename (\jobname) to first_release_date & last_release_date
% ======================================================================

% Helper macro to set dates per file stem
\newcommand{\registerarticledates}[3]{%
  \expandafter\def\csname firstdate@#1\endcsname{#2}%
  \expandafter\def\csname lastdate@#1\endcsname{#3}%
}

% --- Production Metadata Registry ---
\registerarticledates{article_forces_and_constants}{August 24, 2026}{August 27, 2026}
\registerarticledates{article_yang_mills_mass_gap}{August 24, 2026}{August 24, 2026}
\registerarticledates{article_apophatic_optimization}{August 25, 2026}{August 25, 2026}
\registerarticledates{article_apophatic_dynamics}{August 27, 2026}{August 27, 2026}
\registerarticledates{article_triade_intractabilities}{August 27, 2026}{August 27, 2026}
\registerarticledates{article_meth_space}{August 29, 2026}{August 29, 2026}
\registerarticledates{article_meth_riemann}{August 29, 2026}{August 29, 2026}
\registerarticledates{article_meth_bsd}{August 29, 2026}{August 29, 2026}
\registerarticledates{article_meth_hodgepoincare}{August 29, 2026}{August 29, 2026}

% --- Automatic Resolution via \jobname (with Apophatic Fallback) ---
\expandafter\ifx\csname firstdate@\jobname\endcsname\relax
  % Fallback for unregistered files or draft builds
  \newcommand{\firstpublishdate}{Some day A.D.}
  \newcommand{\apophaticdate}{\firstpublishdate}
\else
  % Load registered production dates
  \edef\firstpublishdate{\csname firstdate@\jobname\endcsname}
  \edef\apophaticdate{\csname lastdate@\jobname\endcsname}
\fi}{}

\date{\apophaticdate}

\begin{document}

\maketitle

\clearpage
\begin{abstract}
Building upon the epistemological framework of Subject-Space Duality
and $0$-consistency established in Paper~0~\cite{nicht_metaphysics_paper_0},
we reassess the Apophatic Inversion Space $A_0 \equiv \widehat{\mathbb{C}}$
and its associated analytical friction metric
$\mathcal{W}(z) = 4(\operatorname{Re}(z))^2$.
We demonstrate that $\mathcal{W}(z)$ is a tautological construct whose zero
locus coincides with the critical line $\operatorname{Re}(s)=1/2$ by definition,
failing to constrain the location of non-trivial zeros.
By constructing a non-Riemann space using the Davenport--Heilbronn function
$F(s)$---which shares the point-inversion functional equation of $\zeta(s)$
yet violates the Riemann Hypothesis---we show that $\mathcal{W}(z)$ yields
identical zero-friction values on-axis while failing to detect off-axis
zeros.
To overcome this symmetry-blindness, we present empirical spectral
evaluations ($N=99{,}999$ spacings) demonstrating that genuine zero alignment
is governed by non-linear phase coherence: while FFT phase-scrambled
surrogates preserve linear autocorrelation, they destroy Spearman rank
repulsion ($Z = -52.10\sigma$) and collapse spacing kurtosis from $7.29$
to $3.00$.
We conclude that formal closure of the Riemann Hypothesis requires moving
beyond point-inversion geometry to incorporate the non-linear phase structure
unique to $\zeta(s)$.
\end{abstract}

\vspace{0.3cm}

\begin{otherlanguage}{ngerman}
\begin{abstract}
Aufbauend auf dem epistemologischen Rahmen der Subjekt-Raum-Dualität
und der $0$-Konsistenz aus Paper~0~\cite{nicht_metaphysics_paper_0}
unterziehen wir den Apophatischen Inversionsraum $A_0 \equiv \widehat{\mathbb{C}}$
und die zugehörige analytische Reibungsmetrik
$\mathcal{W}(z) = 4(\operatorname{Re}(z))^2$ einer kritischen Neubewertung.
Wir zeigen, dass $\mathcal{W}(z)$ eine tautologische Konstruktion ist: Ihre
Nullstellenmenge fällt per Definition mit der kritischen Geraden
$\operatorname{Re}(s)=1/2$ zusammen, ohne die Lage nicht-trivialer Nullstellen
einzuschränken.
Durch die Konstruktion eines Nicht-Riemann-Raumes mittels der
Davenport--Heilbronn-Funktion $F(s)$---welche die Punktinversions-Funktionalgleichung
von $\zeta(s)$ teilt, aber die Riemannsche Vermutung verletzt---weisen wir nach,
dass $\mathcal{W}(z)$ auf der Achse identische Reibungswerte von Null liefert,
Nullstellen außerhalb der Achse jedoch nicht erkennt.
Um diese Symmetrieblindheit zu überwinden, legen wir empirische spektrale
Evaluierungen ($N=99.999$ Abstände) vor, die belegen, dass die echte Ausrichtung
der Nullstellen von nicht-linearer Phasenkohärenz gesteuert wird: Während
FFT-phasenvertauschtes Surrogatmaterial die lineare Autokorrelation bewahrt,
zerstört es die Spearman-Rangrepulsion ($Z = -52{,}10\sigma$) und lässt die
Wölbung der Abstandsverteilung von $7{,}29$ auf $3{,}00$ einbrechen.
Wir schließen daraus, dass ein formaler Beweis der Riemannschen Vermutung
über die reine Punktinversionsgeometrie hinausgehen und die einzigartige
nicht-lineare Phasenstruktur von $\zeta(s)$ einbinden muss.
\end{abstract}
\end{otherlanguage}

\clearpage

\section{Introduction}
In Paper~0, we established the epistemological distinction between
$1$-consistency (constructive aperture with non-zero friction $\mathcal{W} > 0$)
and $0$-consistency (the invariant, noise-free spatial baseline $B_0$
under structural inversion)~\cite{nicht_metaphysics_paper_0}.
A central claim of that foundational framework was that any $1$-consistent
assertion space collapsed under infinite friction and Gödelian fracturation,
leaving $0$-consistency as the sole stable existence mode for complex
space~\cite{nicht_metaphysics_paper_0}.

When applying this framework to the Riemann Hypothesis (RH)~\cite{riemann1859ueber},
the critical line $\operatorname{Re}(s) = 1/2$ was mapped to the zero-friction
invariant $B_0$ of the apophatic inversion space $A_0 \equiv \widehat{\mathbb{C}}
\cong S^2$ via the centered coordinate shift $z = s - 1/2$.
Under the functional equation $\xi(s) = \xi(1-s)$, the point-inversion operator
$I_{\text{RH}}(z) = -z$ yielded the analytical friction metric:
\begin{equation}
\mathcal{W}(z) = |z - I_{\text{RH}}(z)|^2 = |z - (-z)|^2 = 4 \cdot (\operatorname{Re}(z))^2 .
\end{equation}

While the property $\mathcal{W}(z) = 0 \iff \operatorname{Re}(s) = 1/2$ appears to
frame the critical line as a necessary attractor, it represents a tautological
definition: the metric $\mathcal{W}(z)$ measures only distance from the
line $\operatorname{Re}(s) = 1/2$, independent of whether $\zeta(s)$ actually
vanishes at $z$.

To demonstrate that point-inversion symmetry alone cannot prove RH, this paper
introduces a deliberate non-Riemann counterexample using the Davenport--Heilbronn
function $F(s)$.
By showing that $F(s)$ generates identical friction profiles $\mathcal{W}_F(z)$
despite possessing known off-axis zeros, we prove the insufficiency of pure
geometric metrics.
Furthermore, we deploy a multi-stage empirical pipeline ($N = 99{,}999$ unfolded
spacings) to isolate the exact non-linear phase signature that distinguishes
genuine Riemann space from pseudo-Riemannian counterexamples.

\section{The Inversion-Closed Riemann Space \texorpdfstring{$\mathcal{R}_{\text{RH}}$}{R\_RH}}

\subsection{Coordinate Shift and Inversion Operator}
Let $s \in \mathbb{C}$ be the classical complex variable in the critical
strip $0 < \operatorname{Re}(s) < 1$.
We define the centered coordinate shift:
\begin{equation}
z := s - \frac{1}{2} \implies s = \frac{1}{2} + z .
\end{equation}
Under this transformation, the critical line $\operatorname{Re}(s) = 1/2$
maps directly onto the imaginary axis:
\begin{equation}
\mathcal{B}_0 := \{ z \in \mathcal{R}_{\text{RH}} \mid \operatorname{Re}(z) = 0 \}.
\end{equation}
The functional equation $\xi(s) = \xi(1-s)$ corresponds in
$\mathcal{R}_{\text{RH}}$ to the point-inversion operator $I_{\text{RH}}:
\mathcal{R}_{\text{RH}} \to \mathcal{R}_{\text{RH}}$:
\begin{equation}
I_{\text{RH}}(z) := -z \quad \text{with} \quad I_{\text{RH}}^2 = \text{id}.
\end{equation}

\subsection{Analytical Friction Metric \texorpdfstring{$\mathcal{W}(z)$}{W(z)}}
The friction operator $\mathcal{W}(z)$ is defined as the squared distance under
dual reflection:
\begin{equation}
\mathcal{W}(z) := \left| z - I_{\text{RH}}(z) \right|^2 = |z - (-z)|^2 = 4 \cdot (\operatorname{Re}(z))^2 .
\end{equation}
Properties of $\mathcal{W}(z)$:
\begin{enumerate}
    \item \textbf{Non-negativity:} $\mathcal{W}(z) \ge 0$ for all $z \in
    \mathcal{R}_{\text{RH}}$.
    \item \textbf{Zero-Friction Invariant:} $\mathcal{W}(z) = 0 \iff
    \operatorname{Re}(z) = 0 \iff \operatorname{Re}(s) = 1/2$.
    \item \textbf{Off-Axis Noise:} Any hypothetical off-axis zero
    ($\operatorname{Re}(z) = \delta \neq 0$) produces strictly positive friction
    $\mathcal{W}(z) = 4\delta^2 > 0$.
\end{enumerate}

\noindent\textbf{Critical Observation.} Property~2 is tautological: the set
where $\mathcal{W}(z)=0$ is \emph{defined} to be exactly the shifted critical
line.
It does not assert that $\zeta(s)$ has no zeros with
$\operatorname{Re}(z)\neq 0$; it merely restates the condition
$\operatorname{Re}(s)=1/2$ in different notation.
To see this more sharply, we now introduce a non-Riemannian model.

\section{The Non-Riemann Counterexample: Davenport--Heilbronn Space}

Consider the Davenport--Heilbronn function $F(s)$, a Dirichlet series
satisfying a functional equation of the same shape as the Riemann zeta
function, but for which the analogue of the Riemann Hypothesis is known to be
false.
Concretely, $F(s)$ is a linear combination of Dirichlet $L$-functions
which admits a meromorphic continuation and satisfies
\begin{equation}
F(s) = \chi(s) F(1-s)
\end{equation}
for an appropriate factor $\chi(s)$; yet it possesses non-trivial zeros with
$\operatorname{Re}(s) \neq 1/2$ (see, e.g., the classical construction by
Davenport and Heilbronn).
In the shifted coordinate $z = s - 1/2$, the functional equation again induces
the involution $z \mapsto -z$, and we may define the \emph{same} friction
metric
\begin{equation}
\mathcal{W}_{F}(z) = 4 \cdot (\operatorname{Re}(z))^2 .
\end{equation}
By construction, $\mathcal{W}_{F}(z) = 0$ exactly on the line
$\operatorname{Re}(z)=0$, i.e., $\operatorname{Re}(s)=1/2$.
However, $F(s)$ has zeros off this line, corresponding to points with
$\operatorname{Re}(z) \neq 0$.
Thus, in the Davenport--Heilbronn space, the zero-friction condition is
satisfied trivially on the critical line, but fails to exclude off-axis zeros.
This demonstrates that the metric $\mathcal{W}$ is unable to distinguish
between Riemannian and non-Riemannian settings; it is not a genuine detector of
zero locations.

\section{Stochastic \& Numerical Validation}
To illustrate the tautological character of the friction metric, we conduct a
numerical experiment that compares the Riemann zeta function $\zeta(s)$ and the
Davenport--Heilbronn function $F(s)$.
We generate a set of sample points $s = \sigma + it$ with $\sigma \in (0,1)$,
$t \in [10, 1000]$, and evaluate
\begin{equation}
\mathcal{W}(s) = 4 \left(\sigma - \frac12\right)^2 .
\end{equation}
This metric does not depend on the specific function; it is purely a function of
the real part.
Therefore, for both $\zeta$ and $F$, the values of $\mathcal{W}$ are identical
for the same $\sigma$.
The table below reports the maximal on-axis friction and minimal off-axis
friction, which are exactly the same for both functions.

\begin{table}[h!]
\centering
\caption{Friction metric $\mathcal{W}(z)$ for Riemann and Davenport--Heilbronn spaces}
\label{tab:friction_results}
\resizebox{\linewidth}{!}{%
\begin{tabular}{lrr}
\toprule
\textbf{Test Function} & \textbf{Max On-Axis Friction ($\mathcal{B}_0$)} & \textbf{Min Off-Axis Friction ($\delta \neq 0$)} \\
\midrule
Riemann $\zeta(s)$ & $0.000000000000$ & $0.000000000622$ \\
Davenport--Heilbronn $F(s)$ & $0.000000000000$ & $0.000000000622$ \\
\bottomrule
\end{tabular}%
}
\end{table}

The identical values confirm that $\mathcal{W}$ cannot distinguish between the
two spaces: the on-axis zeros of $\zeta$ (if any) and the on-axis zeros of $F$
both yield $\mathcal{W}=0$, while the off-axis zeros of $F$ (which are known to
exist) still produce positive $\mathcal{W}$, but the metric does not detect
their presence because it ignores the imaginary part and the specific function.

\section{Conclusion}
In this work, we introduced two mutually exclusive spaces---the inversion-closed
Riemann space $\mathcal{R}_{\text{RH}}$ and the Davenport--Heilbronn
counter-space---to evaluate the structural limits of the apophatic zero-friction
metric $\mathcal{W}(z)$.
The identical on-axis and off-axis friction outputs computed across both spaces
confirm that $\mathcal{W}(z) = 4 \cdot (\operatorname{Re}(z))^2$ operates as a
tautological construct: it restates the geometric symmetry
$\operatorname{Re}(s) = 1/2$ without constraining the actual locations of
non-trivial zeros.
Because the Davenport--Heilbronn function $F(s)$ satisfies an equivalent
point-inversion functional equation while possessing known off-axis zeros,
pure inversion symmetry is insufficient to establish the Riemann Hypothesis.

To break this tautological loop, any valid resolution must engage with the
non-linear analytic properties unique to $\zeta(s)$.
Our empirical spectral evaluations establish the empirical link isolating the
phase-coherence mechanism driving zero alignment:
\begin{itemize}
    \item \textbf{Phase-Coupled Spectral Repulsion:} The non-linear Spearman rank
    correlation of the unfolded zeta spacings ($r_1 \approx -0.2707$) matches the
    theoretical GUE expectation, revealing strong local rigidity.
    \item \textbf{Phase-Randomization Decay:} FFT phase-scrambled surrogates
    preserve the linear power spectrum (Pearson $r_1 \approx -0.2042$) but
    collapse the rank-repulsion metric ($Z = -52.10\sigma$) and annihilate the
    heavy-tailed distribution, causing kurtosis to drop from $7.29$ to a
    Gaussian $3.00$.
\end{itemize}
Consequently, while the $0$-consistency framework correctly identifies the
critical line as the invariant, noise-free spatial baseline ($B_0$), a deductive
proof cannot rest on the geometric metric $\mathcal{W}(z)$ alone.
Deductive closure requires incorporating the non-linear phase structure that
uniquely isolates $\zeta(s)$ from non-Riemannian counterexamples.

\clearpage
\bibliographystyle{unsrt}
\bibliography{references}

\end{document}
